\section{Tổng quan Kiến trúc Hệ thống}
\label{sec:architecture}

Trong lĩnh vực tương tác Người -- Robot hiện đại (\textit{Human-Robot Interaction -- HRI}), việc điều khiển robot di động thông qua các giao diện dòng lệnh hoặc chỉ định tọa độ hình học $(x, y, \theta)$ thuần túy bộc lộ nhiều hạn chế về tính trực quan và khả năng thích ứng của người dùng phổ thông. Nhằm giải quyết thách thức này, bài thực hành tập trung nghiên cứu và hiện thực hóa một hệ thống điều hướng ngữ nghĩa theo thời gian thực (\textit{Autonomous Semantic Navigation}) cho nền tảng robot di động TurtleBot 4 trong môi trường mô phỏng 3D Ignition Gazebo và ROS 2 Humble.

\subsection{Mục tiêu Kỹ thuật và Phạm vi Đề tài}

Hệ thống được thiết kế nhằm đáp ứng trọn vẹn 5 mục tiêu cốt lõi:
\begin{enumerate}[label=\textbf{\arabic*.}]
    \item \textbf{Phân vùng bản đồ tự động (Map Segmentation):} Phân tích bản đồ lưới chiếm chỗ 2D (\textit{Occupancy Grid Map}) kích thước $14\,\text{m} \times 11\,\text{m}$ (độ phân giải $0.05\,\text{m/pixel}$) của khuôn viên \textit{Smart Campus}, tự động trích xuất ranh giới 5 khu vực chức năng gồm: \textbf{Khu A} (Phòng học), \textbf{Khu B} (Phòng thí nghiệm robot), \textbf{Khu C} (Văn phòng hành chính), \textbf{Khu D} (Thư viện) và \textbf{Khu E} (Sảnh chính / Hành lang trung tâm) thông qua thuật toán Watershed có kiểm soát và biến đổi khoảng cách Euclidean.
    \item \textbf{Mô hình hóa Ngữ nghĩa Không gian (Semantic Spatial Modeling):} Gắn kết thông tin định danh, tên gọi đa ngôn ngữ (Việt -- Anh), hệ thống từ đồng nghĩa phong phú (\textit{aliases}), tính toán tâm hình học (\textit{centroid}), đường bao lồi (\textit{convex hull}) và tự động sinh điểm tiếp cận điều hướng an toàn (\textit{goal waypoint}) tại tâm cửa mở rộng $1.3\,\text{m}$.
    \item \textbf{Nhận dạng Giọng nói Tiếng Việt Chống Méo Quá Ngưỡng (Robust Speech-to-Text Pipeline):} Thiết kế chuỗi tiền xử lý âm thanh số bao gồm bộ lọc thông cao 80\,Hz, khử trôi mức một chiều (\textit{DC offset}), nắn méo bão hòa mềm (\textit{$\tanh$ de-clipping}), bộ ngắt khoảng lặng thông minh (\textit{Smart Noise Gate / VAD}) và tự động cân bằng âm lượng thích nghi (\textit{Adaptive Auto AGC}). Tích hợp kiến trúc nhận dạng đa tầng kết hợp giữa \textbf{Google Cloud Deep Learning Speech AI} (độ chính xác $>98\%$) và mô hình \textbf{Vosk Kaldi Offline} có cơ chế lọc nhiễu từ khóa điều hướng (\textit{Domain Keyword Guard}).
    \item \textbf{Phân tích Ngôn ngữ Tự nhiên & Lập kế hoạch Đa chặng (NLU & Multi-hop Planning):} Xây dựng bộ phân tích cú pháp (\textit{Natural Language Understanding -- NLU}) dựa trên văn phạm biểu thức chính quy mở rộng, hỗ trợ nhận dạng ý định di chuyển đơn điểm, tuần tự đa chặng liên phòng (*"tới khu E rồi tới văn phòng"*), dừng khẩn cấp và tra cứu trạng thái.
    \item \textbf{Điều hướng Tự trị Nav2 & Giao diện Studio 2D Thời gian thực:} Kết nối Nav2 Action Client (\textit{NavigateToPose}), tối ưu hóa bộ điều khiển DWB Local Planner loại bỏ hiện tượng đơ dừng và dao động góc tại cửa; đồng thời xây dựng giao diện điều khiển \textit{3D Robotics Studio} bằng PyQt5 hỗ trợ Push-to-Talk, VU meter 3 màu và hiển thị bản đồ trực quan.
\end{enumerate}

\subsection{Sơ đồ Kiến trúc Phân tầng Hệ thống}

Kiến trúc phần mềm tổng thể được tổ chức theo mô hình phân tầng chặt chẽ với các luồng dữ liệu một chiều rõ ràng, không có mũi tên chéo cắt ngang các khối chức năng (Hình~\ref{fig:system_architecture}).

\begin{figure}[H]
\centering
\begin{tikzpicture}[
    node distance=1.1cm and 0.9cm,
    block/.style={draw=uetblue, line width=0.8pt, rounded corners=3pt, fill=white, align=center, font=\small, inner sep=6pt},
    flow/.style={-{Latex[length=2.5mm]}, draw=uetblue, line width=1.0pt}
]
    % Layer 1: Input & Data Sources
    \node[block, fill=softblue, text width=4.6cm] (mic) {\textbf{Tín hiệu Micro (16\,kHz)}\\PCM 32-bit Float Stream};
    \node[block, fill=softgreen, right=1.2cm of mic, text width=4.6cm] (grid) {\textbf{Bản đồ Lưới Chiếm chỗ}\\\code{map.pgm} + \code{map.yaml}};
    \node[block, fill=softpurple, right=1.2cm of grid, text width=4.6cm] (sim) {\textbf{Mô phỏng 3D Gazebo}\\TurtleBot 4 + Smart Campus};

    % Layer 2: Perception & Preprocessing
    \node[block, fill=softblue, below=1.1cm of mic, text width=4.6cm] (cleaner) {\textbf{Bộ Tiền Xử lý Âm thanh}\\HPF 80\,Hz + AGC + Noise Gate};
    \node[block, fill=softgreen, below=1.1cm of grid, text width=4.6cm] (seg) {\textbf{Phân vùng Watershed}\\Distance Transform + Seeds};
    \node[block, fill=softpurple, below=1.1cm of sim, text width=4.6cm] (amcl) {\textbf{Định vị Toàn cục}\\AMCL Particle Filter + LiDAR};

    % Layer 3: Cognition & Semantic Mapping
    \node[block, fill=softblue, below=1.1cm of cleaner, text width=4.6cm] (stt) {\textbf{Hybrid STT Engine}\\Google AI (Chính) / Vosk};
    \node[block, fill=softgreen, below=1.1cm of seg, text width=4.6cm] (sem_map) {\textbf{Cơ sở Dữ liệu Ngữ nghĩa}\\\code{zones.yaml} (5 Phân vùng)};
    \node[block, fill=softorange, below=1.1cm of amcl, text width=4.6cm] (gui) {\textbf{3D Robotics Studio GUI}\\PyQt5, VU Meter, Telemetry};

    % Layer 4: Planning, Decision & Control
    \node[block, fill=softorange, below=1.1cm of stt, text width=4.6cm] (nlu) {\textbf{Bộ Phân tích NLU Parser}\\Trích xuất Ý định \& Lộ trình};
    \node[block, fill=softorange, below=1.1cm of sem_map, text width=4.6cm] (mgr) {\textbf{Voice Navigation Manager}\\Quản lý Hàng đợi \& Trạng thái};
    \node[block, fill=softpurple, below=1.1cm of gui, text width=4.6cm] (nav2) {\textbf{Ngăn xếp Điều hướng Nav2}\\Navfn Planner + DWB Controller};

    % Connections - Vertical Downward Flows
    \draw[flow] (mic) -- (cleaner);
    \draw[flow] (cleaner) -- (stt);
    \draw[flow] (stt) -- (nlu);

    \draw[flow] (grid) -- (seg);
    \draw[flow] (seg) -- (sem_map);

    \draw[flow] (sim) -- (amcl);

    % Horizontal & Inter-layer Flows (Orthogonal routing without cross-overs)
    \draw[flow] (nlu) -- (mgr);
    \draw[flow] (sem_map) -- (mgr);
    \draw[flow] (mgr) -- (nav2);

    % Status & Control loops
    \draw[flow] (nav2) -- (sim);
    \draw[flow] (amcl) -- (gui);
    \draw[flow] (mgr) -- (gui);
\end{tikzpicture}
\caption{Sơ đồ kiến trúc phân tầng hệ thống điều hướng ngữ nghĩa TurtleBot 4 (thiết kế theo luồng dữ liệu chuẩn hóa, không có đường cắt chéo).}
\label{fig:system_architecture}
\end{figure}

Quy trình vận hành khép kín gồm 4 bước mạch lạc:
\begin{itemize}
    \item \textbf{Nhánh Âm thanh (Cột 1):} Tín hiệu micro 16\,kHz được chuẩn hóa, lọc sạch tiếng quạt tản nhiệt, chuyển hóa thành văn bản tiếng Việt qua Google AI và đưa vào bộ phân tích cú pháp NLU.
    \item \textbf{Nhánh Không gian (Cột 2):} Bản đồ lưới Occupancy Grid được xử lý hình thái học và phân vùng Watershed để kiến tạo cơ sở dữ liệu ngữ nghĩa \code{zones.yaml}.
    \item \textbf{Tầng Điều phối (Cột 2, Hàng 4):} Node \code{Voice Navigation Manager} liên kết lộ trình NLU với cơ sở dữ liệu phòng và tọa độ AMCL để gửi mục tiêu đến Nav2.
    \item \textbf{Tầng Chấp hành & Trực quan (Cột 3):} Nav2 DWB Controller sinh lệnh vận tốc \code{/cmd_vel} điều khiển robot trong Gazebo, đồng thời cập nhật toàn bộ trạng thái lên giao diện Studio GUI.
\end{itemize}
